\documentclass[11pt]{article}
\usepackage[margin=1in]{geometry}
\usepackage{amsmath,amssymb}
\usepackage{booktabs}
\usepackage{hyperref}
\usepackage{enumitem}
\usepackage{xcolor}
\usepackage{microtype}

\title{AutoResearch-RL: A Constraint-First Autonomous Research Loop for\
Code-Level Optimization Under Extreme Systems Budgets}
\author{AutoResearch-RL Deep Learning Research Team}
\date{April 2026}

\begin{document}
\maketitle

\begin{abstract}
We present a systems-and-learning analysis of \texttt{AutoResearch-RL}, an autonomous framework that treats model-research iteration as a Markov Decision Process (MDP) over \emph{source code mutations}. Instead of searching only in a fixed hyperparameter box, the agent applies structured diffs to a seed training script, while a deterministic orchestration layer enforces strict feasibility constraints (syntax correctness, artifact capacity estimate, causality policy, and wall-clock execution budget). The architecture combines (i) a PPO-style policy controlling mutation behavior, (ii) a robust patch application mechanism, (iii) a CPU pre-check and dispatch pipeline, and (iv) a power-law SPRT filter for early termination of statistically unpromising runs. We provide a detailed semantic decomposition of the codebase, formalize the objective and reward shaping, analyze reliability and failure surfaces, and propose a practical roadmap toward stronger scientific validity and improved sample efficiency.
\end{abstract}

\section{Introduction}
Most deep learning progress still follows a human-in-the-loop pattern: hypothesize a change, edit code, run experiments, inspect outcomes, and repeat. \texttt{AutoResearch-RL} operationalizes this cycle as a continuous control loop in which an RL agent edits training code directly and receives feedback from execution outcomes and constraint checks.

The repository emphasizes three design commitments:
\begin{enumerate}[leftmargin=1.5em]
    \item \textbf{Code-level search}: actions are textual/structural code patches rather than only scalar knob adjustments.
    \item \textbf{Constraint-first execution}: invalid or non-compliant candidates are filtered before expensive compute.
    \item \textbf{Compute-aware triage}: intermediate telemetry is used to abort runs unlikely to beat current state-of-the-art (SOTA).
\end{enumerate}

This document is written as an arXiv-style internal research paper: comprehensive, implementation-grounded, and focused on both algorithmic and systems semantics.

\section{Problem Formulation}
\subsection{MDP over Program Mutations}
Let $c_t$ denote the current best training program at iteration $t$, and let $h_t$ denote a bounded trajectory memory of prior experiments. The policy receives a state summary
\begin{equation}
    s_t = \Phi(c_t, h_t, d_t),
\end{equation}
where $d_t$ is runtime diagnostics/telemetry (e.g., iteration index, OOM indicators, recent abort patterns).

An action $a_t$ is a structured patch (search/replace diffs) generated by the policy+LLM stack. Applying the patch produces candidate code $\tilde{c}_t$. Evaluation yields result object $r_t$ with status, metrics, and error modes. The environment produces scalar reward $R_t$ and updates memory.

\subsection{Objective}
The optimization target is to minimize validation BPB under strict constraints. Conceptually,
\begin{equation}
\max_\pi \; \mathbb{E}_{\pi}\left[\sum_{t=1}^{T} \gamma^{t-1} R_t\right],
\end{equation}
with reward shaped to encode quality improvements and strict penalties for invalid or non-causal behaviors.

\section{System Architecture}
\subsection{End-to-End Control Path}
One iteration follows:
\begin{enumerate}[leftmargin=1.5em]
    \item Load current best code and contextual state.
    \item Policy computes mutation-control signal and obtains JSON patch proposal.
    \item Patch parser applies AST-assisted / whitespace-robust rewrite.
    \item Causality audit + syntax smoke test + capacity estimate gate execution.
    \item Dispatcher runs candidate in isolated process/container and streams JSON telemetry.
    \item SPRT module monitors learning curve and may issue early abort.
    \item Environment computes reward and updates rolling memory buffer.
    \item If candidate beats SOTA, persist artifact and promote as new best code.
\end{enumerate}

\subsection{Asymmetric Compute Planes}
The design intentionally separates:
\begin{itemize}[leftmargin=1.5em]
    \item \textbf{CPU orchestration plane}: cheap checks, policy context assembly, dispatch control, and logging.
    \item \textbf{GPU evaluation plane}: expensive forward/backward evidence generation for promising candidates only.
\end{itemize}
This asymmetry is central to throughput economics in autonomous search loops.

\section{Component Analysis from the Codebase}
\subsection{Meta-Agent and Diff Semantics}
\texttt{agent/ppo\_agent.py} provides three key layers:
\begin{itemize}[leftmargin=1.5em]
    \item \textbf{State vector extraction}: scalarized context from history and code-derived signals.
    \item \textbf{Policy-value network}: compact actor-critic module producing a control variable used as LLM temperature.
    \item \textbf{Patch pipeline}: parse JSON diff payloads and apply changes via AST-first matching with textual fallback.
\end{itemize}

This architecture is pragmatic: it integrates symbolic editing with a lightweight policy backbone. The trade-off is limited expressivity and potentially weak long-horizon credit assignment under sparse improvements.

\subsection{Environment and Reward Modeling}
\texttt{agent/mdp\_env.py} implements a shaped reward:
\begin{align}
R_t &= \underbrace{\alpha(\mathrm{BPB}_{\text{sota}}-\mathrm{BPB}_t)}_{\text{quality improvement}}
+ \underbrace{\beta \exp(\text{staleness})}_{\text{novelty incentive}}
- \underbrace{P_{\text{syntax/capacity/causality/waste}}}_{\text{constraint penalties}}.
\end{align}

Notably, abort penalties scale with lateness of early stopping, approximating wasted compute.

\subsection{Orchestrator and Dispatch Reliability}
\texttt{orchestrator/orchestrator.py} validates syntax via AST parse and estimates artifact footprint by combining compressed source size with mixed-precision parameter bytes (int6 + bf16 approximations). Candidates that fail basic feasibility checks are rejected prior to dispatch.

\texttt{orchestrator/docker\_runner.py} executes candidate scripts, consumes line-stream JSON telemetry, applies timeout control, and integrates SPRT-triggered termination.

\subsection{SPRT for Online Run Triage}
\texttt{gpu\_cluster/sprt.py} fits a power-law loss model,
\begin{equation}
L(t) = a\, t^{-b} + c,
\end{equation}
and estimates confidence around asymptotic behavior using covariance from nonlinear fit. If projected best-case remains above threshold+margin, run termination is triggered. A plateau heuristic adds robustness against stalled dynamics.

\subsection{Causality Auditor}
\texttt{auditor/causality\_auditor.py} performs AST heuristics to catch suspicious forward-looking indexing/slicing patterns and negative shifts/rolls. The seed script also includes runtime-style causality assertions during attention/evaluation simulation.

\section{Golden Seed: Architectural Prior and Search Surface}
The seed program (\texttt{seed/train\_gpt.py}) defines the mutation substrate. It includes:
\begin{itemize}[leftmargin=1.5em]
    \item Simulated int6 linear modules with per-row scales.
    \item QK-gain normalized attention using causal attention calls.
    \item Dynamic MLP expansion and residual transformer blocks.
    \item Depth recurrence (reusing a small physical block set across loops).
    \item Dynamic low-rank additive path ($\text{LoRA-like}$ loop deltas).
    \item Muon optimizer and EMA-like warmdown averaging.
\end{itemize}

From a search perspective, this is a high-information prior: compact yet structurally rich, allowing policy edits to produce large behavioral shifts from local textual changes.

\section{Safety, Validity, and Failure Surfaces}
\subsection{What Is Well-Engineered}
\begin{itemize}[leftmargin=1.5em]
    \item Explicit guardrails for syntax, capacity, time budget, and causality policy.
    \item Structured telemetry and persistent experiment logs.
    \item Early-stop strategy to preserve expensive compute.
    \item Modular separation that supports local testing of major subsystems.
\end{itemize}

\subsection{Current Risk Envelope}
\begin{itemize}[leftmargin=1.5em]
    \item Several paths are still simulation-oriented (mock patches/telemetry assumptions).
    \item Causality checks are heuristic and may under/over-fire on nontrivial transformations.
    \item Policy learning is lightweight; single-step update flow may underutilize trajectory structure.
    \item Capacity estimates are approximate abstractions, not full bit-exact packaging audits.
\end{itemize}

\section{Experimental Protocol Recommendations}
To mature toward publishable empirical rigor:
\begin{enumerate}[leftmargin=1.5em]
    \item \textbf{Deterministic benchmarking harness}: fixed seeds, fixed workloads, and repeated trials with confidence intervals.
    \item \textbf{Ablation grid}: isolate contribution of novelty bonus, SPRT thresholding, and causality auditing.
    \item \textbf{Patch quality metrics}: acceptance rate, syntactic validity rate, semantic rollback rate, and average improvement latency.
    \item \textbf{Cost accounting}: wall-clock/GPU-hour budget per net SOTA gain.
    \item \textbf{Leakage stress tests}: adversarial mutation suites to evaluate auditor robustness.
\end{enumerate}

\section{Roadmap for Research-Grade Evolution}
\subsection{Short Horizon (Engineering Stability)}
\begin{itemize}[leftmargin=1.5em]
    \item Add schema-level validation for patch payloads and AST-preserving transforms.
    \item Improve dispatcher protocol with explicit heartbeat/step cadence guarantees.
    \item Strengthen failure taxonomy and attach remediation suggestions to each class.
\end{itemize}

\subsection{Mid Horizon (Learning Quality)}
\begin{itemize}[leftmargin=1.5em]
    \item Transition to batched, multi-step PPO updates with generalized advantage estimation.
    \item Learn patch priors from accepted-diff corpora.
    \item Add retrieval over historical experiment motifs for memory-aware proposal generation.
\end{itemize}

\subsection{Long Horizon (Scientific Generalization)}
\begin{itemize}[leftmargin=1.5em]
    \item Expand action space beyond model script into tokenizer and eval policy co-design.
    \item Introduce multi-objective optimization (BPB, latency, robustness, package entropy).
    \item Build formal verification hooks for causal-evaluation invariants.
\end{itemize}

\section{Conclusion}
\texttt{AutoResearch-RL} demonstrates a compelling pattern: combine policy-guided code mutation, strict systems guardrails, and compute-aware run triage to automate meaningful portions of deep learning research iteration. Its current implementation is strongest as a robust prototype architecture with clear pathways to research-grade maturity. The central conceptual contribution is not a single optimizer trick, but a \emph{closed-loop research operating system} where code edits are first-class actions and feasibility constraints are embedded as environment physics.

\section*{Reproducibility Notes}
This paper was produced from repository-level analysis of the current implementation and specifications. It is intentionally grounded in implemented module semantics rather than external benchmark claims.

\end{document}
